% 子文件：9.黎曼积分.tex
% 主文档使用 \input{9.黎曼积分} 引入，不含 document 环境

\section{黎曼积分}
\subsection{黎曼积分的概念}

设$f(P)$是几何形体$\varOmega$上有定义的点函数。将$\varOmega$分割为$n$个小的几何形体，同时用它们表示其度量。称数，记$\lambda$。任取，作乘积的和式
\[
\sum_{i=1}^{n}f(P_i)\Delta\varOmega_i
\]
如果不论分割如何，极限
\[
\lim_{\lambda\rightarrow0}\sum_{i=1}^{n}f(P_i)\Delta\varOmega_i
\]
都存在，且为同一个值，则称此极限值为函数$f(P)$在几何形体$\varOmega$上的黎曼积分，即
\[
\int_\varOmega f(P)\,d\varOmega
\]
此时也说$f(P)$在$\varOmega$上可积。

定理9.1\quad 若

\subsection{黎曼积分的性质}
\begin{enumerate}
    \item 当$f(P)=1$时，
    \item 线性性质
    \item 对积分域的可加性质
    \item 比较性质
    \item 估值性质
    \item 积分中值定理
    \item 对称性质
\end{enumerate}

\subsection{黎曼积分的分类}
\begin{enumerate}
    \item 二重积分

当

    \item 三重积分

当几何形体

    \item 对弧长的曲线积分

当几何形体

    \item 对面积的曲面积分
\end{enumerate}

\section{二重积分的计算}
\subsection{直角坐标系下二重积分的计算}

设$D$为$Oxy$平面上一有界闭域，$f(P)$的二重积分
\[
\iint_D f(x,y)\,dxdy
\]
存在。这时二重积分可表示为。当，即表示，其面积为。由已知，体积。习惯上，记作，并把。得到一个计算公式。公式把二重积分化为累次积分。计算时。

当，不等式组表示，计算公式。公式将二重积分化为另一种累次积分，先。

若函数$f(x,y)$在积分。

\subsection{极坐标系下二重积分的计算}

在极坐标系下，设，
\[
\iint_D f(x,y)\,dxdy=\iint_{D'} f(r\cos\theta,r\sin\theta)\,r\,dr d\theta
\]
其中

\section{三重积分的计算}
\subsection{直角坐标系下三重积分的计算}

在直角坐标系下体积微元是
\[
dV=dx\,dy\,dz
\]

\subsubsection{投影法（先一后二法）}

设，其中。先用，再用，故有公式。这就是。

当，即积分域。由公式。

\subsubsection{截面法（先二后一法）}

\subsection{柱坐标系下三重积分的计算}

设点$P(x,y,z)$在$Oxy$坐标面上的投影。

柱坐标的三组坐标面是：
\begin{itemize}
    \item $r=$常数，
    \item $\theta=$常数，
    \item $z=$常数。
\end{itemize}

显然，直角坐标与柱坐标的关系是
\[
\begin{cases}
x=r\cos\theta,\\
y=r\sin\theta,\\
z=z.
\end{cases}
\]
用三组坐标面分割积分域，故柱坐标下体积微元是
\[
dV=r\,dr d\theta dz
\]
因此柱坐标系下三重积分可写为
\[
\iiint_V f(x,y,z)\,dV
=\iiint_{V'} f(r\cos\theta,r\sin\theta,z)\,r\,dr d\theta dz.
\]

柱坐标系下。从而。故。直角坐标系下的关系是。当。

\subsection{球坐标系下三重积分的计算}

设。

球坐标的三组坐标面是：
\begin{enumerate}
    \item
    \item
    \item
\end{enumerate}

点$P$的直角坐标$(x,y,z)$与球坐标$(\rho,\phi,\theta)$之间的关系是
\[
x=\rho\sin{\phi}\cos{\theta},\quad
y=\rho\sin{\phi}\sin{\theta},\quad
z=\rho\cos{\phi}.
\]
用球坐标的三组坐标面分割积分域$V$，典型的小块是直角六面体。所以，在球坐标系下，体积微元是
\[
dV=\rho^2\sin{\phi}\,d\rho d\phi d\theta
\]
故球坐标系下三重积分可以写为
\[
\iiint_V f(x,y,z)\,dV
=\iiint_{V'} f(\rho\sin\phi\cos\theta,\rho\sin\phi\sin\theta,\rho\cos\phi)\,\rho^2\sin\phi\,d\rho d\phi d\theta.
\]

球坐标，则。直角坐标系下。当。

\section{第一型曲线积分的计算}

设$l$是，由参数方程
\[
\begin{cases}
x=x(t),\\
y=y(t),
\end{cases}\quad \alpha\le t\le\beta
\]
给出。如果函数，则对弧长的曲线积分
\[
\int_l f(x,y)\,ds
\]
存在。

设点。故由定积分定义知。于是对弧长的曲线积分可化为定积分计算，
\[
\int_l f(x,y)\,ds
=\int_{\alpha}^{\beta}f\bigl(x(t),y(t)\bigr)\sqrt{\bigl(x'(t)\bigr)^2+\bigl(y'(t)\bigr)^2}\,dt.
\]
注意，这里弧长微元$ds$就是弧微分，
\[
ds=\sqrt{(dx)^2+(dy)^2}.
\]
当！由此可见。这是第一型曲线积分的一个特性。

如果$l$是：则有公式1。

\section{第一型曲面积分的计算}

设空间曲面$\Sigma$为，其中，则曲面积分
\[
\iint_\Sigma f(x,y,z)\,dS
\]
存在。

如果，显然有，其中，故。从而，其中，定义知。于是。

\section{黎曼积分的应用举例}
\subsection{物体的质心}

由静力学知。

\subsection{转动惯量}

转动惯量。